\section{物理层 I：景观、势垒与类能量直觉}

如果被谨慎使用，物理语言仍然很有用。
类能量直觉之所以有帮助，是因为它压缩了几个定性思想：
\begin{itemize}[nosep]
  \item 深势阱难以逃离；
  \item 浅势阱容易失稳；
  \item 邻近的山脊或鞍点决定了一个小推动能否改变路径；
  \item 即便没有完美控制，噪声也可能导致跨越势垒。
\end{itemize}

\begin{discussion}
景观语言的教学收益，不在于读者真的掌握了某个字面意义上的机械能函数。
收益在于，它为持续与逃逸提供了一个可处理的心智模型。
因此，在本书中，能量图景应当被视为一种受中等车道数学推理支撑、并且受生物学中关于稳定性与转变的强实证正当性支撑的结构化类比。
\end{discussion}